% GENERATED FILE. Edit canonical lessons, then run npm run generate:books.
% canonical-source-hash: fnv1a64:2ae15ca239aa8982
% canonical-lessons: KA-C237-kaccu, KA-C237-balasu, KA-C237-saripadisu, KA-C237-hode, KA-C237-prayatnisu

\chapter{To bite, To use, To repair, To hit, To try}
\label{ch:ka-to-bite-to-use-to-repair-to-hit-to-try}
\hlchaptermodality{237}

\begin{chapteropening}
\textbf{By the end of this chapter:} \emph{I can say to bite, to use, to repair, to hit and to try.}

Complete the last lesson of chapter 237: I can say to bite, to use, to repair, to hit and to try.
\end{chapteropening}

\section[\texorpdfstring{\kn{ಕಚ್ಚು} (kaccu)}{ಕಚ್ಚು (kaccu)}]{\kn{ಕಚ್ಚು} (kaccu) — to bite}

\label{lesson:KA-C237-kaccu}



\begin{quote}
Before the new one: say the Kannada for to bargain, then the Kannada for to climb.
\end{quote}

\subsection*{You'll want to know: \kn{ಕಚ್ಚು}}
\textbf{\kn{ಕಚ್ಚು}} — \emph{kaccu} — \textquotedblleft{}to bite\textquotedblright{}.

\begin{grammarlens}[title={What you've built}]
One more word for this chapter.
\end{grammarlens}

\subsection*{Your turn}
\begin{itemize}
\raggedright
  \item \emph{Say it:} \emph{kaccu}
  \item \emph{Say it:} \emph{kaccu}, once more
  \item \emph{Say it:} say \emph{kaccu}
  \item \emph{Recall:} say the Kannada for to bargain, then the Kannada for to climb, then say \emph{kaccu} again
\end{itemize}

\begin{tcolorbox}[breakable,colback=teal!4,colframe=teal!35!black,title={Before you move on}]
What does \kn{ಕಚ್ಚು} mean? (\textquotedblleft{}To bite\textquotedblright{}.) Say it once more.
\end{tcolorbox}
\section[\texorpdfstring{\kn{ಬಳಸು} (baḷasu)}{ಬಳಸು (baḷasu)}]{\kn{ಬಳಸು} (baḷasu) — to use}

\label{lesson:KA-C237-balasu}



\begin{quote}
Before the new one: say the Kannada for to climb, then the Kannada for to bite.
\end{quote}

\subsection*{You'll want to know: \kn{ಬಳಸು}}
\textbf{\kn{ಬಳಸು}} — \emph{baḷasu} — \textquotedblleft{}to use\textquotedblright{}.

\begin{grammarlens}[title={What you've built}]
One more word for this chapter.
\end{grammarlens}

\subsection*{Your turn}
\begin{itemize}
\raggedright
  \item \emph{Say it:} \emph{baḷasu}
  \item \emph{Say it:} \emph{baḷasu}, once more
  \item \emph{Say it:} say \emph{baḷasu}
  \item \emph{Recall:} say the Kannada for to climb, then the Kannada for to bite, then say \emph{baḷasu} again
\end{itemize}

\begin{tcolorbox}[breakable,colback=teal!4,colframe=teal!35!black,title={Before you move on}]
What does \kn{ಬಳಸು} mean? (\textquotedblleft{}To use\textquotedblright{}.) Say it once more.
\end{tcolorbox}
\section[\texorpdfstring{\kn{ಸರಿಪಡಿಸು} (saripaḍisu)}{ಸರಿಪಡಿಸು (saripaḍisu)}]{\kn{ಸರಿಪಡಿಸು} (saripaḍisu) — to repair, to fix}

\label{lesson:KA-C237-saripadisu}



\begin{quote}
Before the new one: say the Kannada for to bite, then the Kannada for to use.
\end{quote}

\subsection*{You'll want to know: \kn{ಸರಿಪಡಿಸು}}
\textbf{\kn{ಸರಿಪಡಿಸು}} — \emph{saripaḍisu} — \textquotedblleft{}to repair, to fix\textquotedblright{}.

\begin{grammarlens}[title={What you've built}]
One more word for this chapter.
\end{grammarlens}

\subsection*{Your turn}
\begin{itemize}
\raggedright
  \item \emph{Say it:} \emph{saripaḍisu}
  \item \emph{Say it:} \emph{saripaḍisu}, once more
  \item \emph{Say it:} say \emph{saripaḍisu}
  \item \emph{Recall:} say the Kannada for to bite, then the Kannada for to use, then say \emph{saripaḍisu} again
\end{itemize}

\begin{tcolorbox}[breakable,colback=teal!4,colframe=teal!35!black,title={Before you move on}]
What does \kn{ಸರಿಪಡಿಸು} mean? (\textquotedblleft{}To repair, to fix\textquotedblright{}.) Say it once more.
\end{tcolorbox}
\section[\texorpdfstring{\kn{ಹೊಡೆ} (hoḍe)}{ಹೊಡೆ (hoḍe)}]{\kn{ಹೊಡೆ} (hoḍe) — to hit, to beat}

\label{lesson:KA-C237-hode}



\begin{quote}
Before the new one: say the Kannada for to use, then the Kannada for to repair.
\end{quote}

\subsection*{You'll want to know: \kn{ಹೊಡೆ}}
\textbf{\kn{ಹೊಡೆ}} — \emph{hoḍe} — \textquotedblleft{}to hit, to beat\textquotedblright{}.

\begin{grammarlens}[title={What you've built}]
One more word for this chapter.
\end{grammarlens}

\subsection*{Your turn}
\begin{itemize}
\raggedright
  \item \emph{Say it:} \emph{hoḍe}
  \item \emph{Say it:} \emph{hoḍe}, once more
  \item \emph{Say it:} say \emph{hoḍe}
  \item \emph{Recall:} say the Kannada for to use, then the Kannada for to repair, then say \emph{hoḍe} again
\end{itemize}

\begin{tcolorbox}[breakable,colback=teal!4,colframe=teal!35!black,title={Before you move on}]
What does \kn{ಹೊಡೆ} mean? (\textquotedblleft{}To hit, to beat\textquotedblright{}.) Say it once more.
\end{tcolorbox}
\section[\texorpdfstring{\kn{ಪ್ರಯತ್ನಿಸು} (prayatnisu)}{ಪ್ರಯತ್ನಿಸು (prayatnisu)}]{\kn{ಪ್ರಯತ್ನಿಸು} (prayatnisu) — to try}

\label{lesson:KA-C237-prayatnisu}



\begin{quote}
Before the new one: say the Kannada for to repair, then the Kannada for to hit.
\end{quote}

\subsection*{You'll want to know: \kn{ಪ್ರಯತ್ನಿಸು}}
\textbf{\kn{ಪ್ರಯತ್ನಿಸು}} — \emph{prayatnisu} — \textquotedblleft{}to try\textquotedblright{}.

\begin{grammarlens}[title={What you've built}]
One more word for this chapter.
\end{grammarlens}

\subsection*{Your turn}
\begin{itemize}
\raggedright
  \item \emph{Say it:} \emph{prayatnisu}
  \item \emph{Say it:} \emph{prayatnisu}, once more
  \item \emph{Say it:} say \emph{prayatnisu}
  \item \emph{Recall:} say the Kannada for to repair, then the Kannada for to hit, then say \emph{prayatnisu} again
\end{itemize}

\begin{tcolorbox}[breakable,colback=teal!4,colframe=teal!35!black,title={Before you move on}]
What does \kn{ಪ್ರಯತ್ನಿಸು} mean? (\textquotedblleft{}To try\textquotedblright{}.) Say it once more.
\end{tcolorbox}
